\documentclass[10pt]{article}
\usepackage[margin=0.7in]{geometry}
\usepackage[T1]{fontenc}
\usepackage[utf8]{inputenc}
\usepackage{lmodern}
\usepackage{microtype}
\usepackage{array}
\usepackage{booktabs}
\usepackage{longtable}
\usepackage{tabularx}
\usepackage{pdflscape}
\usepackage[hidelinks]{hyperref}
\usepackage{fancyhdr}
\usepackage{titlesec}

\setlength{\headheight}{14pt}
\pagestyle{fancy}
\fancyhf{}
\lhead{Open-Source Deployment Orchestration}
\rhead{Comparison Matrix}
\cfoot{\thepage}
\titleformat{\section}{\large\bfseries}{\thesection}{0.6em}{}
\setlength{\parindent}{0pt}
\setlength{\parskip}{0.5em}
\emergencystretch=2em
\urlstyle{same}
\renewcommand{\arraystretch}{1.2}

\title{Open-Source Deployment Orchestration\\\large Comparison Matrix by Officially Documented Scope}
\author{}
\date{May 2026}

\begin{document}
\maketitle

\textbf{Reading note.} This matrix compares tools by their documented primary scope, operating model, and best-fit use case. It is intentionally taxonomy-first: tools are grouped by the type of orchestration they are designed to perform, not by broad marketing labels.\cite{KubernetesOverview,ArgoCDHome,AnsiblePlaybooks,AirflowHome,TemporalWhatIs,KubeflowIntro}

\begin{landscape}
\small
\begin{longtable}{>{\raggedright\arraybackslash}p{2.3cm} >{\raggedright\arraybackslash}p{2.2cm} >{\raggedright\arraybackslash}p{3.6cm} >{\raggedright\arraybackslash}p{3.3cm} >{\raggedright\arraybackslash}p{2.5cm} >{\raggedright\arraybackslash}p{3.0cm} >{\raggedright\arraybackslash}p{3.2cm}}
\toprule
\textbf{Tool} & \textbf{Category} & \textbf{Primary documented scope} & \textbf{Execution model} & \textbf{Infra assumption} & \textbf{Best fit} & \textbf{Key boundary} \\
\midrule
\endfirsthead
\toprule
\textbf{Tool} & \textbf{Category} & \textbf{Primary documented scope} & \textbf{Execution model} & \textbf{Infra assumption} & \textbf{Best fit} & \textbf{Key boundary} \\
\midrule
\endhead
\bottomrule
\endfoot
\textbf{Kubernetes} & Platform orchestrator & Manages containerized workloads and services with declarative configuration and automation.\cite{KubernetesOverview} & Cluster scheduler and control plane for containerized applications.\cite{KubernetesOverview} & Kubernetes cluster required & General-purpose container orchestration at scale & Not a GitOps controller or data-pipeline engine \\

\textbf{Argo CD} & GitOps / CD & Continuously compares live Kubernetes state to desired state in Git and syncs drift.\cite{ArgoCDHome,ArgoCDDeclarative} & Kubernetes controller that reconciles desired and live state.\cite{ArgoCDHome} & Kubernetes required & Git-driven deployment and reconciliation for Kubernetes apps & Does not replace the runtime orchestrator \\

\textbf{Nomad} & Platform orchestrator & Deploys and manages containerized and legacy applications using a unified workflow.\cite{NomadWhatIs} & Workload scheduler for containers, legacy apps, microservices, and batch jobs.\cite{NomadWhatIs} & Nomad cluster required & Mixed workload estates including non-containerized apps & Not a Kubernetes GitOps controller \\

\textbf{Docker Swarm} & Platform orchestrator & Manages a cluster of Docker daemons and deploys services across swarm nodes.\cite{DockerSwarm,DockerSwarmConcepts} & Docker service orchestration across swarm managers and workers.\cite{DockerSwarm,DockerSwarmConcepts} & Docker swarm mode required & Simpler Docker-native clustered deployments & Narrower scope and ecosystem than Kubernetes-style platforms \\

\textbf{Ansible} & Config orchestration & Manages configurations and deployments; sequences multi-tier rollouts and delegated actions.\cite{AnsiblePlaybooks,AnsibleRolling} & Playbook-driven remote execution and orchestration.\cite{AnsiblePlaybooks} & Managed hosts and inventory required & Host automation, environment setup, ordered rollouts & Not a workload scheduler \\

\textbf{Apache Airflow} & Data workflow & Develops, schedules, and monitors batch-oriented workflows represented as DAGs.\cite{AirflowHome,AirflowArchitecture} & DAG-based task orchestration for batch workflows.\cite{AirflowArchitecture} & Airflow deployment required & Scheduled data pipelines and batch dependencies & Not designed as a durable business-process runtime \\

\textbf{Prefect} & Data workflow & Builds, deploys, runs, and monitors data pipelines while tracking dependencies and failures.\cite{PrefectQuickstart} & Flow-based orchestration of Python-defined workflows.\cite{PrefectQuickstart} & Prefect open-source or cloud-backed setup & Python-native data workflows with operational controls & Not a cluster scheduler \\

\textbf{Dagster} & Data workflow & Data orchestration with lineage, observability, declarative programming, and testability.\cite{DagsterOverview} & Asset- and job-oriented orchestration for data systems.\cite{DagsterOverview} & Dagster deployment required & Data platforms where asset awareness and lineage matter & Not primarily a generic application orchestrator \\

\textbf{Temporal} & Durable workflows & Provides durable workflow execution that preserves state and progress across failures.\cite{TemporalWhatIs,TemporalWorkflowExecution} & Durable workflow runtime for long-running application code.\cite{TemporalWhatIs} & Temporal service required & Stateful business processes, retries, timers, and compensation logic & Not primarily a batch DAG scheduler \\

\textbf{Flyte} & ML / AI platform & Orchestrates AI workflows and positions itself as a Kubernetes-native workflow orchestrator.\cite{FlyteOSS} & Workflow orchestration platform oriented to AI and data workloads.\cite{FlyteOSS} & Flyte platform; commonly Kubernetes-centered & AI workflow development and productionization & More specialized than general CD or host automation tools \\

\textbf{Kubeflow} & ML / AI platform & Provides a foundation of tools for AI platforms on Kubernetes across the AI lifecycle.\cite{KubeflowIntro} & Modular Kubernetes-native AI platform components.\cite{KubeflowIntro,KubeflowPipelines} & Kubernetes required & Building an AI platform stack on Kubernetes & Broader AI platform scope than a standalone scheduler \\

\textbf{Metaflow} & ML / AI platform & Human-friendly Python library for developing, deploying, and operating data-intensive apps, especially data science, ML, and AI.\cite{MetaflowWhatIs,MetaflowWhy} & Python-centric workflow model for data-intensive applications.\cite{MetaflowWhatIs} & Metaflow runtime environment required & Teams prioritizing Pythonic ML and data-science workflows & Not a generic infrastructure orchestrator \\
\end{longtable}
\normalsize
\end{landscape}

\section{Interpretation guidance}
\begin{itemize}
    \item Compare tools first by \emph{category}, then by \emph{execution model}. Most bad selections come from collapsing those two concepts into one.\cite{KubernetesOverview,ArgoCDHome,TemporalWhatIs}
    \item When two tools appear adjacent in the matrix, that does not mean they are substitutes. For example, Kubernetes and Argo CD are commonly complementary, not competing.\cite{KubernetesOverview,ArgoCDHome}
    \item Airflow, Prefect, Dagster, and Temporal all orchestrate work, but they do so for meaningfully different runtime models.\cite{AirflowArchitecture,PrefectQuickstart,DagsterOverview,TemporalWhatIs}
\end{itemize}

\begin{thebibliography}{99}
\bibitem{KubernetesOverview} Kubernetes Documentation. \emph{Overview}. Available: \url{https://kubernetes.io/docs/concepts/overview/}
\bibitem{ArgoCDHome} Argo CD Documentation. \emph{Argo CD - Declarative GitOps CD for Kubernetes}. Available: \url{https://argo-cd.readthedocs.io/en/stable/}
\bibitem{ArgoCDDeclarative} Argo CD Documentation. \emph{Declarative Setup}. Available: \url{https://argo-cd.readthedocs.io/en/stable/operator-manual/declarative-setup/}
\bibitem{NomadWhatIs} HashiCorp Nomad Documentation. \emph{What is Nomad?}. Available: \url{https://developer.hashicorp.com/nomad/docs/what-is-nomad}
\bibitem{DockerSwarm} Docker Docs. \emph{Swarm mode}. Available: \url{https://docs.docker.com/engine/swarm/}
\bibitem{DockerSwarmConcepts} Docker Docs. \emph{Swarm mode key concepts}. Available: \url{https://docs.docker.com/engine/swarm/key-concepts/}
\bibitem{AnsiblePlaybooks} Ansible Community Documentation. \emph{Working with playbooks}. Available: \url{https://docs.ansible.com/projects/ansible/latest/playbook_guide/playbooks.html}
\bibitem{AnsibleRolling} Ansible Community Documentation. \emph{Playbook Example: Continuous Delivery and Rolling Upgrades}. Available: \url{https://docs.ansible.com/projects/ansible/latest/playbook_guide/guide_rolling_upgrade.html}
\bibitem{AirflowHome} Apache Airflow Documentation. \emph{What is Airflow?}. Available: \url{https://airflow.apache.org/docs/apache-airflow/stable/index.html}
\bibitem{AirflowArchitecture} Apache Airflow Documentation. \emph{Architecture Overview}. Available: \url{https://airflow.apache.org/docs/apache-airflow/stable/core-concepts/overview.html}
\bibitem{PrefectQuickstart} Prefect Documentation. \emph{Quickstart}. Available: \url{https://docs.prefect.io/v3/get-started/quickstart}
\bibitem{DagsterOverview} Dagster Documentation. \emph{Overview}. Available: \url{https://docs.dagster.io/}
\bibitem{TemporalWhatIs} Temporal Documentation. \emph{What is Temporal?}. Available: \url{https://docs.temporal.io/temporal}
\bibitem{TemporalWorkflowExecution} Temporal Documentation. \emph{Workflow Execution overview}. Available: \url{https://docs.temporal.io/workflow-execution}
\bibitem{FlyteOSS} Flyte OSS Documentation. \emph{User guide / Flyte OSS}. Available: \url{https://docs.flyte.org/}
\bibitem{KubeflowIntro} Kubeflow Documentation. \emph{Introduction}. Available: \url{https://www.kubeflow.org/docs/started/introduction/}
\bibitem{KubeflowPipelines} Kubeflow Documentation. \emph{Kubeflow Pipelines Overview}. Available: \url{https://www.kubeflow.org/docs/components/pipelines/overview/}
\bibitem{MetaflowWhatIs} Metaflow Documentation. \emph{What is Metaflow}. Available: \url{https://docs.metaflow.org/introduction/what-is-metaflow}
\bibitem{MetaflowWhy} Metaflow Documentation. \emph{Why Metaflow}. Available: \url{https://docs.metaflow.org/introduction/why-metaflow}
\end{thebibliography}

\end{document}
